% Generated draft: requires mathematical and source review.
\documentclass[11pt]{article}
\usepackage[margin=25mm]{geometry}
\usepackage{fontspec}
\setmainfont{Latin Modern Roman}
\usepackage{amsmath,amssymb,mathtools,unicode-math}
\setmathfont{Latin Modern Math}
\usepackage{enumitem,needspace,xurl,hyperref}
\setlength{\parindent}{0pt}
\setlength{\parskip}{5pt}
\setlength{\emergencystretch}{3em}
\newcommand{\N}{\mathbb{N}}
\newcommand{\Z}{\mathbb{Z}}
\newcommand{\Q}{\mathbb{Q}}
\newcommand{\R}{\mathbb{R}}
\newcommand{\C}{\mathbb{C}}
\newcommand{\F}{\mathbb{F}}
\newcommand{\A}{\mathbb{A}}
\newcommand{\PP}{\mathbb{P}}
\newcommand{\E}{\mathbb{E}}
\newcommand{\eps}{\varepsilon}
\newcommand{\dd}{\,\mathrm{d}}
\begin{document}
\section*{UnsolvedMath ID: 30004093}
\textbf{Source problem code:} catalogue entry
\subsection{Definitions and mathematical statement}
Let \(p\) be a prime and let \(K\) be the \(p\)-adic field appearing in the supplied record. Write \(\mathcal O_K\) for its valuation ring, \(\mathfrak m_K\) for the maximal ideal, and \(k_K=\mathcal O_K/\mathfrak m_K\) for the residue field. The record does not explicitly specify whether \(K\) must be a finite extension of \(\Q_p\).

Under the customary integral, origin-fixing interpretation of a Lubin pair, the power series belong to \(T\mathcal O_K[[T]]\). Such a series \(g\) is called stable when \(g'(0)\ne0\) and \(g'(0)\) is not a root of unity. Invertibility here means invertibility for composition in \(T\mathcal O_K[[T]]\), equivalently \(g'(0)\in\mathcal O_K^\times\). These conventions are needed to interpret the statement; the supplied record does not itself define ``stable'' or ``Lubin pair.'' With these conventions, the hypotheses are
\[
 f,u\in T\mathcal O_K[[T]],\qquad f\circ u=u\circ f,
\]
\[
 0\ne f'(0)\in\mathfrak m_K,\qquad
 u'(0)\in\mathcal O_K^\times\text{ is not a root of unity},\qquad
 \overline f\ne0\text{ in }k_K[[T]].
\]
The last condition is coefficientwise reduction modulo \(\mathfrak m_K\).

For clarity, a one-dimensional commutative formal group law over a commutative ring \(R\) is a series \(S(X,Y)\in R[[X,Y]]\) satisfying
\[
 S(X,0)=X,\quad S(0,Y)=Y,\quad S(X,Y)=S(Y,X),
\]
\[
 S(S(X,Y),Z)=S(X,S(Y,Z)),
\]
and admitting a formal inverse. An endomorphism of \(S\) is a series \(a\in TR[[T]]\) satisfying
\[
 a(S(X,Y))=S(a(X),a(Y)).
\]
This is the usual formal-group setting suggested by the question; the record does not expressly specify dimension, commutativity, or the coefficient ring of \(S\).

\textbf{Question, subject to the unresolved qualifications below.} For every pair \((f,u)\) satisfying the supplied Lubin-pair hypotheses, does there exist a formal group \(S\) and endomorphisms \(a,b\) of \(S\) to which \(f,u\) are simultaneously semiconjugate?

To make this conclusion precise, a common coefficient ring and an admissible nonzero series \(h\), with \(h(0)=0\), must be specified. One possible direction of simultaneous semiconjugacy is
\[
 h\circ f=a\circ h,\qquad h\circ u=b\circ h;
\]
the reverse direction is
\[
 f\circ h=h\circ a,\qquad u\circ h=h\circ b.
\]
The supplied record does not determine the direction, whether the same \(h\) is required for both series, the permitted coefficient extensions, or additional conditions on \(h\). These alternatives are not asserted to be equivalent. Consequently, this is a qualified statement-only formulation, and a fully specified conjecture requires source clarification.
\subsection{Short English statement}
The supplied problem asks whether every commuting stable Lubin pair, consisting of a noninvertible series \(f\) with nonzero reduction and an invertible series \(u\), admits a simultaneous formal-group model by semiconjugacy. The integral coefficient convention, the formal group's base ring, and the precise semiconjugacy requirements need clarification. The supplied status is partially solved; that literature assessment has not been independently verified.
\subsection{Sources}
\begin{itemize}
\item Supplied local UnsolvedMath registry record, permanent ID \(30004093\), title: Formal Group Models for Commuting p-Adic Power Series. Exact source problem code: \texttt{catalogue entry}. Collection: Oberwolfach Reports Open Problems. No individual proposer is identified in the supplied record.
\item Supplied report citation: \emph{Non-Archimedean Geometry and Applications (2020)}, Oberwolfach Reports, DOI \texttt{10.4171/owr/2019/8}, pp.\ 513--575. \url{https://doi.org/10.4171/owr/2019/8}. The record separately gives proposed year 2019 and the tag \texttt{owr-year:2019}; these are retained without reconciling them with the citation's 2020 designation.
\item Supplied literature reference, labeled ``Berger on Lubin's conjecture'': \url{https://arxiv.org/abs/1702.06037}.
\item Supplied literature reference, labeled ``Height-one Lubin conjecture'': \url{https://arxiv.org/abs/2607.03257}.
\item Supplied literature reference, labeled ``Oberwolfach report'': \url{https://doi.org/10.4171/owr/2019/8}.
\item Literature assessment checked at \texttt{2026-08-22}: this is supplied dated curation, not an independently verified assessment. It labels the problem partially solved and states that Lubin's conjecture is proved for substantial families and height-one systems, while the general assertion that every relevant commuting \(p\)-adic dynamical system is semiconjugate to a formal-group system remains open. The supplied background calls this dated literature triage derived from OWR curation and explicitly requests independent verification of time-sensitive status and cited claims.
\item The supplied statement status is \texttt{corrected\_verified}. Its accompanying qualification says that mathematical content was preserved while notation, typography, and sentence structure were repaired for publication. This records the supplied editorial assessment, not independent mathematical verification of this draft.
\end{itemize}
\end{document}
